\documentclass[noauthor,noinstructornotes]{ximera}

\title{Week 13}
\author{Math 1193 Design Team}

\begin{document}
\begin{abstract}
    Activities for Week 13.
\end{abstract}
\maketitle

\begin{problem}
Let's play a game! Divide into pairs and choose who will be
Player 1 and Player 2. Below are the instructions.

\begin{enumerate}
    \item Player 1 starts with a hidden equation. On the left-hand side
    is a number, and on the right-hand side is something complicated
    with an \(x\) in it. For example:
    \begin{itemize}
        \item \(2 = 3x - 9\)
    \end{itemize}
    \item Player 1 then writes down on a scrap of paper each step/operation
    applied to the \(x\) on the right-hand side, keeping in mind PEMDAS.
    For example:
    \begin{itemize}
        \item Multiply by 3, subtract 9
    \end{itemize}
    \item Player 1 gives the paper to Player 2 without showing the original
    equation.
    \item Player 2 then looks at the paper and writes down the steps
    that need to be applied to both sides of the equation to undo
    the process described by Player 1. For example:
    \begin{itemize}
        \item Add 9, divide by 3
    \end{itemize}
    \item Both players then take Player 2's list and apply it to both sides
    of the equation. What is the result of this process?
\end{enumerate}

Play the game with the following equations:
\begin{enumerate}
    \item \(1 = 2(x - 5)\)
    \item \(4 = \log(x)\)
    \item \(5 = \ln(x + 3) - 1\)
    \item \(2 = 5 \cdot 9^x\) 
\end{enumerate}
\end{problem}

\begin{problem}
Julia worked on solving the equation \(\ln(x) + \ln(x - 21) = 2\). Her
work is as follows:

\begin{align*}
\ln(x) + \ln(x - 21) & = 2 & \text{(original equation)}\\
e^{\ln(x)} + e^{ln(x - 21)} & = e^2 & \text{(exponentiating on both sides)} \\
x + (x - 21) & = e^2 & \text{(canceling logs and exponents)}\\
2x - 21 & = e^2 & \text{(combining like terms)}\\
2x & = e^2 + 21 & \text{(adding 21 on both sides)}\\
x & = \frac{e^2 + 21}{2} & \text{(dividing by 2 on both sides)}\\
\end{align*}

Is Julia's work correct? How can you tell? If it's not correct, identify
the step in which she made a mistake and try to find a solution.
\end{problem}

\end{document}
